% Options for packages loaded elsewhere
\PassOptionsToPackage{unicode}{hyperref}
\PassOptionsToPackage{hyphens}{url}
\PassOptionsToPackage{dvipsnames,svgnames,x11names}{xcolor}
\documentclass[
  english,
  10pt,
  a4paper,
]{article}
\usepackage{xcolor}
\usepackage[top=22mm,bottom=20mm,left=18mm,right=18mm]{geometry}
\usepackage{amsmath,amssymb}
\setcounter{secnumdepth}{-\maxdimen} % remove section numbering
\usepackage{iftex}
\ifPDFTeX
  \usepackage[T1]{fontenc}
  \usepackage[utf8]{inputenc}
  \usepackage{textcomp} % provide euro and other symbols
\else % if luatex or xetex
  \usepackage{unicode-math} % this also loads fontspec
  \defaultfontfeatures{Scale=MatchLowercase}
  \defaultfontfeatures[\rmfamily]{Ligatures=TeX,Scale=1}
\fi
\usepackage{lmodern}
\ifPDFTeX\else
  % xetex/luatex font selection
\fi
% Use upquote if available, for straight quotes in verbatim environments
\IfFileExists{upquote.sty}{\usepackage{upquote}}{}
\IfFileExists{microtype.sty}{% use microtype if available
  \usepackage[]{microtype}
  \UseMicrotypeSet[protrusion]{basicmath} % disable protrusion for tt fonts
}{}
\makeatletter
\@ifundefined{KOMAClassName}{% if non-KOMA class
  \IfFileExists{parskip.sty}{%
    \usepackage{parskip}
  }{% else
    \setlength{\parindent}{0pt}
    \setlength{\parskip}{6pt plus 2pt minus 1pt}}
}{% if KOMA class
  \KOMAoptions{parskip=half}}
\makeatother
\usepackage{longtable,booktabs,array}
\usepackage{caption}
\captionsetup[table]{skip=6pt}
\newcounter{none} % for unnumbered tables
\usepackage{calc} % for calculating minipage widths
% Correct order of tables after \paragraph or \subparagraph
\usepackage{etoolbox}
\makeatletter
\patchcmd\longtable{\par}{\if@noskipsec\mbox{}\fi\par}{}{}
\makeatother
\ifLuaTeX
\usepackage[bidi=basic,shorthands=off]{babel}
\else
\usepackage[bidi=default,shorthands=off]{babel}
\fi
\ifLuaTeX
  \usepackage{selnolig} % disable illegal ligatures
\fi
\setlength{\emergencystretch}{3em} % prevent overfull lines
\providecommand{\tightlist}{%
  \setlength{\itemsep}{0pt}\setlength{\parskip}{0pt}}
% ---- Çılgın Kurye document theme: same look as the CSE 396 project proposal ----
% Compile with XeLaTeX (also works on Overleaf). Fonts: TeX Gyre Termes (Times), DejaVu Sans Mono.
\usepackage{fontspec}
\setmainfont{texgyretermes}[Extension=.otf, UprightFont=*-regular, BoldFont=*-bold, ItalicFont=*-italic, BoldItalicFont=*-bolditalic]
\setsansfont{texgyreheros}[Extension=.otf, UprightFont=*-regular, BoldFont=*-bold, ItalicFont=*-italic, BoldItalicFont=*-bolditalic]
\setmonofont{DejaVuSansMono}[Extension=.ttf, UprightFont=*, BoldFont=*-Bold, ItalicFont=*-Oblique, BoldItalicFont=*-BoldOblique, Scale=0.80]
\setmathfont{texgyretermes-math.otf}

\usepackage{xcolor}
\definecolor{ckblue}{HTML}{1F3A5F}
\definecolor{cklight}{HTML}{E8EEF5}
\definecolor{ckgrey}{HTML}{F2F2F2}
\definecolor{ckline}{HTML}{555555}
\definecolor{shadecolor}{HTML}{F4F5F7}

% Headings as in the proposal: serif, bold; subsections bold italic. Numbers are part of the heading text.
\setcounter{secnumdepth}{0}
\usepackage{titlesec}
\titleformat{\section}{\rmfamily\Large\bfseries}{}{0pt}{}
\titleformat{\subsection}{\rmfamily\large\bfseries\itshape}{}{0pt}{}
\titleformat{\subsubsection}{\rmfamily\normalsize\bfseries}{}{0pt}{}
\titlespacing*{\section}{0pt}{12pt}{4pt}
\titlespacing*{\subsection}{0pt}{8pt}{3pt}

% Header and footer (the first page has only the footer, like the proposal)
\usepackage{fancyhdr,lastpage}
\pagestyle{fancy}
\fancyhf{}
\fancyhead[L]{\footnotesize\itshape CSE 396 · Group 8 · Çılgın Kurye}
\fancyhead[R]{\footnotesize\itshape\ckmodule}
\fancyfoot[C]{\footnotesize Page \thepage\ of \pageref*{LastPage}}
\renewcommand{\headrulewidth}{0.4pt}
\setlength{\headheight}{14pt}
\fancypagestyle{ckfirst}{\fancyhf{}\renewcommand{\headrulewidth}{0pt}\fancyfoot[C]{\footnotesize Page \thepage\ of \pageref*{LastPage}}}

% Inline code: breaks only after . : ( , / _ and spaces, never inside a name
\usepackage[obeyspaces,spaces]{url}
\AtBeginDocument{%
  \renewcommand{\UrlBreaks}{\do\.\do\:\do\(\do\,\do\/\do\_\do\|}%
  \renewcommand{\UrlBigBreaks}{}%
  \renewcommand{\UrlNoBreaks}{}%
  \urlstyle{tt}}

% Tables: boxed grid as in the proposal, smaller text
\usepackage{etoolbox}
\AtBeginEnvironment{longtable}{\footnotesize\renewcommand{\arraystretch}{1.2}}
\setlength{\LTpre}{6pt}
\setlength{\LTpost}{8pt}
\setlength{\arrayrulewidth}{0.4pt}

% Code blocks: long lines wrap; diagrams stay in one piece
\usepackage{framed}
\usepackage{fvextra}
\RecustomVerbatimEnvironment{verbatim}{Verbatim}{fontsize=\footnotesize, breaklines}
\makeatletter
\@ifundefined{Shaded}{}{\renewenvironment{Shaded}{\begin{snugshade}}{\end{snugshade}}}
\@ifundefined{Highlighting}{}{\RecustomVerbatimEnvironment{Highlighting}{Verbatim}{commandchars=\\\{\}, fontsize=\footnotesize, breaklines, breakanywhere}}
\makeatother

% Figures drawn with TikZ
\usepackage{tikz}
\usetikzlibrary{arrows.meta,positioning,fit,calc,backgrounds,shapes.geometric}
\tikzset{
  ck/.style={font=\footnotesize, >={Stealth[length=5pt]}, line width=0.6pt},
  box/.style={draw=ckline, fill=white, rounded corners=2pt, align=center, inner sep=4pt, minimum height=7mm},
  mod/.style={box, fill=ckgrey},
  this/.style={box, fill=cklight, draw=ckblue, line width=0.9pt},
  group/.style={draw=ckline, dashed, rounded corners=3pt, inner sep=6pt},
  lbl/.style={font=\scriptsize, align=center, inner sep=1.5pt},
  msg/.style={->, draw=ckline},
}
\usepackage{caption}
\captionsetup[figure]{font=small, labelfont=bf, skip=4pt}
\newcommand{\ckfigure}[2]{\begin{figure}[htbp]\centering #1\caption{#2}\end{figure}}

\usepackage{enumitem}
\setlist{itemsep=1pt, topsep=3pt}
\setlength{\parskip}{4pt}
\setlength{\parindent}{0pt}
\raggedbottom
% ---- end of theme ----
\newcommand{\ckmodule}{M1 Game Server and Rules}
\usepackage{bookmark}
\IfFileExists{xurl.sty}{\usepackage{xurl}}{} % add URL line breaks if available
\urlstyle{same}
% fallback for those not using the hyperref driver hyperxmp:
\makeatletter
\@ifundefined{xmpquote}{\newcommand{\xmpquote}[1]{#1}}{}
\makeatother
\hypersetup{
  pdftitle={M1 Game Server and Rules},
  pdflang={en},
  colorlinks=true,
  linkcolor={black},
  filecolor={Maroon},
  citecolor={Blue},
  urlcolor={black},
  pdfcreator={LaTeX}}

\title{M1 Game Server and Rules --- Module Documentation}
\author{}
\date{}


% ---- colours and drawings ----
% ---- colours ----
\definecolor{srv}{HTML}{2E6FB7}\colorlet{srvL}{srv!15}
\definecolor{cli}{HTML}{2E9E5B}\colorlet{cliL}{cli!15}
\definecolor{adm}{HTML}{E08A1E}\colorlet{admL}{adm!18}
\definecolor{brd}{HTML}{7B4FB3}\colorlet{brdL}{brd!15}
\definecolor{tcpc}{HTML}{1F5FA8}
\definecolor{udpc}{HTML}{D9531E}
\definecolor{obs}{HTML}{D7263D}
\definecolor{tA}{HTML}{00A6D6}
\definecolor{tB}{HTML}{3BB273}
\definecolor{tC}{HTML}{F2C14E}
\definecolor{tD}{HTML}{9B5DE5}

% ---- styles ----
\tikzset{
  nb/.style={rounded corners=3pt, line width=0.9pt, align=center, inner sep=4pt, minimum height=9mm, font=\footnotesize},
  sv/.style={nb, draw=srv, fill=srvL},
  cl/.style={nb, draw=cli, fill=cliL},
  ad/.style={nb, draw=adm, fill=admL},
  bd/.style={nb, draw=brd, fill=brdL},
  hwb/.style={nb, draw=black!50, fill=black!5},
  tcpa/.style={draw=tcpc, line width=1.5pt},
  udpa/.style={draw=udpc, line width=1.5pt, dash pattern=on 5pt off 2.5pt},
  inna/.style={draw=black!55, line width=1pt},
  al/.style={font=\scriptsize, fill=white, inner sep=1.2pt, rounded corners=1pt},
  grp/.style={draw=black!35, dashed, rounded corners=5pt, inner sep=7pt},
  glbl/.style={font=\scriptsize\bfseries, text=black!65},
  hl/.style={line width=2.6pt},
  st/.style={nb, rounded corners=8pt, minimum width=24mm},
  co/.style={draw=black!60, line width=0.5pt},
  cot/.style={font=\scriptsize, align=left, text=black!80},
}

% ---- system map; argument = module to highlight (M1 ... M8, or none) ----
\newcommand{\systemmap}[1]{%
\tikzset{M1/.style={},M2/.style={},M3/.style={},M4/.style={},M5/.style={},M6/.style={},M7/.style={},M8/.style={},none/.style={}}%
\tikzset{#1/.style={hl}}%
\begin{tikzpicture}[ck]
\node[cl, M5, minimum width=27mm] (m5) at (0,3) {\textbf{M5}\\Player phones};
\node[ad, M6, minimum width=27mm] (m6) at (0,0.4) {\textbf{M6}\\Admin panel};
\node[sv, M2, minimum width=27mm, minimum height=38mm] (m2) at (4.6,1.7) {\textbf{M2}\\Network\\and sync};
\node[sv, M1, minimum width=27mm, minimum height=38mm] (m1) at (8.4,1.7) {\textbf{M1}\\Game server\\and rules};
\node[sv, M3, minimum width=25mm] (m3) at (12,2.8) {\textbf{M3}\\Physics};
\node[sv, M4, minimum width=25mm] (m4) at (12,0.6) {\textbf{M4}\\Traffic and AI};
\node[grp, fit=(m2)(m1)(m3)(m4)] (pc) {};
\node[glbl, anchor=south west] at (pc.north west) {MAIN COMPUTER};
\node[bd, M7, minimum width=27mm] (m7) at (4.6,-2.7) {\textbf{M7}\\ESP32 software};
\node[bd, M8, minimum width=27mm] (m8) at (8.4,-2.7) {\textbf{M8}\\LED engine};
\node[hwb, minimum width=25mm] (led) at (12,-2.7) {344 LEDs};
\node[hwb, minimum width=27mm] (sm) at (4.6,-4.4) {Sensor and\\4 motors};
\node[grp, fit=(m7)(m8)(led)(sm)] (bb) {};
\node[glbl, anchor=north east] at (bb.south east) {LED CITY BOARD};
\draw[tcpa, <->] (m5.east) -- node[al]{WebSocket} (m5.east -| m2.west);
\draw[tcpa, <->] (m6.east) -- node[al]{WebSocket} (m6.east -| m2.west);
\draw[tcpa, <->] ([xshift=-6mm]m7.north) -- node[al]{TCP} ([xshift=-6mm]m2.south);
\draw[udpa, <->] ([xshift=6mm]m7.north) -- node[al]{UDP} ([xshift=6mm]m2.south);
\draw[inna, <->] (m2) -- (m1);
\draw[inna, <->] (m1.east |- m3) -- (m3);
\draw[inna, <->] (m1.east |- m4) -- (m4);
\draw[inna, ->] (m7) -- node[al]{state} (m8);
\draw[inna, ->] (m8) -- (led);
\draw[inna, <->] (m7) -- (sm);
\draw[tcpa] (-1.3,-6.0) -- ++(0.9,0) node[right, font=\scriptsize]{TCP: every message arrives, in order};
\draw[udpa] (4.0,-6.0) -- ++(0.9,0) node[right, font=\scriptsize]{UDP: fast, the newest message wins};
\draw[inna] (9.3,-6.0) -- ++(0.9,0) node[right, font=\scriptsize]{inside one program or chip};
\end{tikzpicture}}

% ---- LED board seen from above (6 x 6 roads, 1 unit = 10 cm) ----
\newcommand{\boardbase}{%
\fill[black!88, rounded corners=3pt] (-0.5,-0.5) rectangle (5.5,5.5);
\foreach \i in {0,...,4}{\foreach \j in {0,...,4}{%
  \fill[black!60] (\i+0.13,\j+0.13) rectangle (\i+0.87,\j+0.87);
  \fill[black!22] (\i+0.24,\j+0.24) rectangle (\i+0.52,\j+0.52);
  \fill[black!30] (\i+0.56,\j+0.30) rectangle (\i+0.78,\j+0.62);
  \fill[black!38] (\i+0.28,\j+0.58) rectangle (\i+0.66,\j+0.78);}}
\foreach \i in {0,...,5}{%
  \draw[black, line width=4.5pt] (\i,-0.05) -- (\i,5.05);
  \draw[black, line width=4.5pt] (-0.05,\i) -- (5.05,\i);}
\foreach \i in {0,...,5}{\foreach \k in {0,...,30}{%
  \fill[black!55] (\i,\k/6) circle (0.025);
  \fill[black!55] (\k/6,\i) circle (0.025);}}
\foreach \x/\y/\c in {0/4/tA, 4/4/tB, 0/0/tC, 4/0/tD}{%
  \fill[\c!30!black!60] (\x+0.13,\y+0.13) rectangle (\x+0.87,\y+0.87);
  \fill[\c] (\x+0.3,\y+0.3) rectangle (\x+0.7,\y+0.7);
  \fill[white] (\x+0.4,\y+0.5) circle (0.05); \fill[white] (\x+0.6,\y+0.5) circle (0.05);}}
\newcommand{\nmark}[3]{\node[circle, fill=yellow!85!orange, draw=black, line width=0.4pt, inner sep=0.6pt, minimum size=3.6mm, font=\tiny\bfseries] at (#1,#2) {#3};}
\newcommand{\bike}[4]{%
\fill[#3!35] ($(#1,#2)+(#4:0.40)$) circle (0.05);
\fill[#3!65] ($(#1,#2)+(#4:0.21)$) circle (0.075);
\fill[#3, draw=white, line width=0.4pt] (#1,#2) circle (0.12);}
\newcommand{\dpoint}[3]{\draw[#3!50, line width=0.8pt] (#1,#2) circle (0.26); \draw[#3, line width=1.4pt] (#1,#2) circle (0.16); \fill[#3] (#1,#2) circle (0.07);}
\newcommand{\heavy}[2]{\draw[obs!60!black, line width=4pt] (#1) -- (#2); \draw[obs, line width=4pt, dash pattern=on 3pt off 5pt] (#1) -- (#2);}
\newcommand{\roadwork}[2]{\draw[obs, line width=4pt] (#1) -- (#2);}
\newcommand{\accident}[2]{\draw[obs, line width=4pt] (#1) -- (#2); \draw[white, line width=4pt, dash pattern=on 2.5pt off 3.5pt] (#1) -- (#2);}
\newcommand{\puddle}[2]{\draw[cyan!35, line width=4pt] (#1) -- (#2); \draw[white, line width=1.6pt, dash pattern=on 1pt off 2.5pt] (#1) -- (#2);}
\newcommand{\findme}[2]{\fill[white, opacity=0.25] (#1,#2) circle (0.36); \fill[white, opacity=0.5] (#1,#2) circle (0.24); \fill[white] (#1,#2) circle (0.13);}

% ---- phone game screen (7 x 3.4 units) ----
\newcommand{\phonemock}{%
\fill[black!85, rounded corners=9pt] (-0.3,-0.3) rectangle (7.3,3.7);
\fill[cyan!20] (0,1.7) rectangle (7,3.4);
\fill[cli!14] (0,0) rectangle (7,1.7);
\filldraw[fill=srv!30, draw=srv!70!black] (0.2,1.5) rectangle (1.0,3.0);
\filldraw[fill=adm!35, draw=adm!70!black] (1.1,1.6) rectangle (1.8,2.4);
\filldraw[fill=brd!30, draw=brd!70!black] (1.9,1.65) -- (2.6,1.65) -- (2.6,2.2) -- (2.25,2.5) -- (1.9,2.2) -- cycle;
\filldraw[fill=adm!35, draw=adm!70!black] (4.5,1.65) rectangle (5.1,2.3);
\filldraw[fill=srv!30, draw=srv!70!black] (5.2,1.55) rectangle (6.0,3.1);
\filldraw[fill=cli!35, draw=cli!70!black] (6.1,1.6) rectangle (6.85,2.6);
\fill[brown!60!black] (4.17,1.6) rectangle (4.23,1.8); \fill[cli!70!black] (4.2,1.95) circle (0.18);
\fill[black!55] (1.8,0) -- (5.2,0) -- (3.75,1.7) -- (3.25,1.7) -- cycle;
\draw[white, line width=0.8pt, dash pattern=on 4pt off 4pt] (3.5,0) -- (3.5,1.7);
\filldraw[fill=black!15, draw=black!70] (3.62,1.12) rectangle (3.88,1.32);
\fill[obs] (3.05,1.22) -- (3.15,1.48) -- (3.25,1.22) -- cycle;
\fill[black!85] (3.43,0.08) rectangle (3.57,0.3);
\fill[tA] (3.3,0.28) rectangle (3.7,0.6);
\fill[brown!65] (3.34,0.6) rectangle (3.66,0.82);
\fill[tA!60!black] (3.5,0.95) circle (0.13);
\node[fill=black!70, text=white, rounded corners=2pt, font=\scriptsize\bfseries, inner sep=2pt] at (3.5,3.15) {2:41};
\foreach \c/\s/\x in {tA/120/0.35,tB/95/0.95,tC/140/1.55,tD/80/2.15}
  \node[fill=\c, text=white, font=\tiny\bfseries, rounded corners=1pt, inner sep=1.5pt, minimum width=5mm] at (\x,3.15) {\s};
\fill[white, opacity=0.9] (6.45,2.85) circle (0.38); \draw[black!60] (6.45,2.85) circle (0.38);
\draw[-{Stealth[length=5pt]}, line width=1.8pt, obs] (6.45,2.85) ++(230:0.2) -- ++(50:0.5);
\fill[white, opacity=0.85] (0.55,0.5) circle (0.32); \fill[black!70] (0.43,0.5) -- (0.65,0.63) -- (0.65,0.37) -- cycle;
\fill[white, opacity=0.85] (1.4,0.5) circle (0.32); \fill[black!70] (1.52,0.5) -- (1.3,0.63) -- (1.3,0.37) -- cycle;
\node[fill=obs!80, text=white, rounded corners=3pt, font=\tiny\bfseries, minimum width=8mm, minimum height=6mm] at (5.6,0.5) {Brake};
\node[fill=cli, text=white, rounded corners=3pt, font=\tiny\bfseries, minimum width=8mm, minimum height=8mm] at (6.5,0.6) {Gas};
\node[fill=tC, text=black, rounded corners=3pt, font=\tiny\bfseries, minimum width=9mm, minimum height=5mm] at (6.15,1.45) {Deliver};
\node[fill=white, text=black, rounded corners=3pt, font=\tiny\bfseries, minimum width=9mm, minimum height=5mm] at (0.9,1.45) {Find me};}

% ---- side cut of the board (14 units wide) ----
\newcommand{\boardside}{%
\fill[brown!20] (-0.6,-0.35) rectangle (14.6,0);
\foreach \x/\h/\c in {1.0/0.7/black!18, 2.2/1.1/black!22, 4.6/1.9/black!25, 5.8/2.6/black!28, 7.4/2.2/black!25, 8.6/1.6/black!22, 10.9/1.0/black!18}
  \filldraw[fill=\c, draw=black!45] (\x,1.5) rectangle (\x+0.9,1.5+\h);
\filldraw[fill=tB!35, draw=tB!70!black] (12.4,1.5) rectangle (13.4,2.3);
\fill[black!70, rounded corners=1pt] (0.5,0) rectangle (1.1,0.9);
\fill[black!70, rounded corners=1pt] (12.9,0) rectangle (13.5,0.9);
\filldraw[fill=brown!45, draw=brown!70!black] (0,0.9) rectangle (14,1.3);
\fill[black!80] (0,1.3) rectangle (14,1.38);
\foreach \x in {0.3,0.9,...,13.8} \fill[black!35] (\x,1.42) circle (0.06);
\foreach \x/\c in {3.3/tA, 3.9/tA, 6.9/obs, 7.5/obs, 10.5/tC} \fill[\c] (\x,1.42) circle (0.07);
\fill[white, opacity=0.55] (0,1.38) rectangle (14,1.52);
\filldraw[fill=brd!40, draw=brd] (2.3,0.35) rectangle (4.3,0.85);
\node[font=\tiny] at (3.3,0.6) {ESP32};
\filldraw[fill=brd, draw=brd!60!black] (6.75,0.62) rectangle (7.25,0.88);
\filldraw[fill=adm!60, draw=adm!70!black] (1.8,0.6) circle (0.24);
\filldraw[fill=adm!60, draw=adm!70!black] (12.2,0.6) circle (0.24);
\filldraw[fill=black!15, draw=black!50] (9.3,0.35) rectangle (10.9,0.85);
\node[font=\tiny] at (10.1,0.6) {Power};}

\begin{document}

\thispagestyle{ckfirst}
\begin{center}
{\bfseries GEBZE TECHNICAL UNIVERSITY\par}
{Computer Engineering Department\par}
{CSE 396 Module Documentation\par}
\vspace{8pt}
{\LARGE\bfseries Çılgın Kurye\par}
\vspace{3pt}
{\large\itshape Module M1: Game Server and Rules\par}
\vspace{6pt}
\textbf{Instructor:} Erkan Zergeroğlu\hspace{1.4em}\textbf{Company:} Board\&Byte\hspace{1.4em}\textbf{Group:} Group-8\hspace{1.4em}\textbf{Date:} 6 November 2026
\end{center}
\begin{center}\footnotesize\renewcommand{\arraystretch}{1.15}
\begin{tabular}{|p{0.27\linewidth}|p{0.15\linewidth}|p{0.27\linewidth}|p{0.15\linewidth}|}\hline
\textbf{Name} & \textbf{Student No.} & \textbf{Name} & \textbf{Student No.} \\ \hline
Tuğberk Akbulut & 230104004070 & Abdullatif Yılmaz & 230104004031 \\ \hline
Berk Andaçgil & 230104004066 & Selim Barut & 220104004046 \\ \hline
Umutcan Ocak & 220104004903 &  &  \\ \hline
\end{tabular}
\end{center}
\vspace{2pt}



% ---- table and list helpers ----
\newcolumntype{L}[1]{>{\raggedright\arraybackslash}p{\dimexpr#1\linewidth-2\tabcolsep-\arrayrulewidth\relax}}
\newcolumntype{C}[1]{>{\centering\arraybackslash}p{\dimexpr#1\linewidth-2\tabcolsep-\arrayrulewidth\relax}}
\newcommand{\ckdone}{\item[$\boxtimes$]}
\newcommand{\ckopen}{\item[$\square$]}
\newcommand{\mdot}{$\bullet$}
\newcounter{msec}
\newcounter{chap}
% ---- end of helpers ----
\newcommand{\msection}[1]{\stepcounter{msec}\section{\arabic{msec}. #1}}

\msection{Purpose}

M1 is the brain of the game. It runs on the main computer and it is the only part that can change
the game. It keeps the game clock and the one true game state. It knows every team, store, package,
delivery point, score, fall, obstacle and earthquake. Phones only send their buttons. M1 decides
what happens, and then everybody sees the same result.

\begin{longtable}[]{|L{0.18}|L{0.82}|}
\hline Runs on & Main computer, in one server program together with M2, M3 and M4 \\ \hline
Technology & TypeScript on Node.js; fixed-step game loop at 30 ticks per second; seeded random numbers; JSON files for the city plan and the game log \\ \hline
Requirements & R1 R2 R4 R5 R6 R7 R9 R11 \\ \hline
\end{longtable}

\begin{figure}[htbp]\centering
\systemmap{M1}
\caption{Where M1 sits in the system. The thick box is M1.}
\end{figure}

\msection{Requirements covered}

\begin{longtable}[]{|L{0.07}|L{0.25}|L{0.68}|}
\hline \textbf{ID} & \textbf{Requirement} & \textbf{What M1 does for it} \\ \hline \endhead
R1 & Multiplayer game & Lobby with 4 teams and at most 2 players in a team; bots fill empty places; a game is at most 5 minutes; the team with the most points wins \\ \hline
R2 & Delivery cycle & Automatic pickup at the store; picks a delivery point by travel time; near or far chance per player \\ \hline
R4 & Collisions and falls & Gives a penalty time after every fall that M3 reports \\ \hline
R5 & No map on the phone & Puts the compass target and the ``Find me'' flag into the game state \\ \hline
R6 & Obstacles & Keeps the obstacle list that M3, M4, M5 and M8 use \\ \hline
R7 & Earthquake & One path for a hand hit, the motors and the admin button; road work and falls inside the radius \\ \hline
R9 & Live admin panel & Applies admin commands and sends each parameter to the module that uses it \\ \hline
R11 & Single authority & Only M1 changes the game state; every tick can be played again from the log \\ \hline
\end{longtable}

\msection{Inputs and outputs}

\begin{longtable}[]{|L{0.08}|L{0.20}|L{0.18}|L{0.54}|}
\hline \textbf{Dir.} & \textbf{Data} & \textbf{From / to} & \textbf{Content} \\ \hline \endhead
In & Controls & M2, from phones & Throttle and brake 0--1, steer −1 to 1, deliver; the newest one in each tick \\ \hline
In & Join, ``Find me'' & M2, from phones & Player key, name, team; ``Find me'' request \\ \hline
In & Admin commands & M2, from admin & Parameter changes, start, pause, end, obstacles, earthquake \\ \hline
In & Board events & M2, from board & Earthquake strength 0--10, centre, zone, confidence \\ \hline
In & Traffic & M4 & NPC cars, bot controls, travel-time table \\ \hline
In & Physics & M3 & New motorcycle positions, collision and fall events \\ \hline
Out & Game state & M2, to phones and admin & Every tick: phase, time, scores, players, NPC cars, obstacles, earthquake \\ \hline
Out & Board state & M2, to board & Players, delivery points, obstacles, earthquake in board coordinates \\ \hline
Out & Metrics & M2, to admin & Twice per second: the values of all admin graphs \\ \hline
Out & Game log & Disk & The inputs of every tick, one JSON line per tick \\ \hline
\end{longtable}

\msection{Main components}

\begin{longtable}[]{|L{0.20}|L{0.52}|L{0.28}|}
\hline \textbf{Component} & \textbf{What it does} & \textbf{Related modules} \\ \hline \endhead
Game loop & Runs the steps of a tick in a fixed order, 30 times per second & M3, M4 \\ \hline
Match and lobby & Game phases, teams, bots, keeps a place for 30 s after a player drops & M2, M4 \\ \hline
Delivery and scoring & Pickup, delivery check with a reason when it says no, choice of the delivery point, points & M4, M5, M8 \\ \hline
Penalties & Penalty time after a fall; asks M3 to stand the motorcycle up again & M3 \\ \hline
Obstacles and earthquake & Obstacle list with a version number; earthquake strength, radius, road work and falls & M3, M4, M7 \\ \hline
Admin and parameters & Checks commands, keeps every parameter with its limits, sends each group to its owner & M2, M3, M4, M6 \\ \hline
State builder and recorder & Builds the frozen game state, the board state and the metrics; writes and replays the game log & M2 \\ \hline
\end{longtable}

\msection{How it works}

\textbf{The tick.} The game moves in small steps called ticks. There are 30 ticks in one second, so
one tick has about 33 ms. A tick always runs the same eight steps in the same order, as the figure
below shows. First M1 takes the inputs that M2 has collected since the last tick. If a phone sent
more than one control message, only the newest one is used. Then M1 applies admin commands and
board events, updates the game phase, asks M4 for traffic and bot controls, and asks M3 to move all
motorcycles. After that the game rules run. At the end M1 builds a frozen copy of the state for M2
and writes one line to the game log. Because the order never changes and all random numbers come
from one seed, the same inputs always give the same game.

\begin{figure}[htbp]\centering
\begin{tikzpicture}[ck, node distance=3.5mm]
\node[cl, minimum width=17mm] (a) {Inputs};
\node[ad, minimum width=17mm, right=of a] (b) {Admin and\\board events};
\node[sv, minimum width=17mm, right=of b] (c) {Game\\phase};
\node[sv, minimum width=17mm, right=of c, fill=tB!20, draw=tB!80!black] (d) {Traffic\\M4};
\node[sv, minimum width=17mm, right=of d, fill=tD!18, draw=tD!80!black] (e) {Physics\\M3};
\node[sv, minimum width=17mm, right=of e, fill=obs!12, draw=obs] (f) {Rules};
\node[sv, minimum width=17mm, right=of f] (g) {Frozen\\state};
\node[hwb, minimum width=17mm, right=of g] (h) {Log};
\foreach \u/\v in {a/b,b/c,c/d,d/e,e/f,f/g,g/h} \draw[inna,->] (\u) -- (\v);
\draw[inna,->, srv] (h.south) -- ++(0,-0.6) -| node[al, pos=0.25]{next tick, 33 ms later} (a.south);
\draw[tcpa,->] (g.north) -- ++(0,0.6) node[above, font=\scriptsize]{to M2};
\end{tikzpicture}
\caption{The eight steps of one tick.}
\end{figure}

\textbf{Game flow.} A game has five phases. In the lobby players join and choose a team. When the
admin presses start, a 3 second countdown begins and nobody can move. Then the game runs for at most
5 minutes. The admin can pause and resume the game at any time. When the time ends, M1 shows the
final scores and the game is finished. If a phone loses its connection, M1 keeps that player's
place for 30 seconds. If the player does not come back, a bot takes the motorcycle and continues.

\begin{figure}[htbp]\centering
\begin{tikzpicture}[ck]
\node[st, draw=cli, fill=cliL] (l) at (0,0) {Lobby};
\node[st, draw=adm, fill=admL] (c) at (3.6,0) {Countdown};
\node[st, draw=srv, fill=srvL, line width=1.4pt] (r) at (7.2,0) {Running};
\node[st, draw=brd, fill=brdL] (f) at (10.8,0) {Finished};
\node[st, draw=black!50, fill=black!6] (p) at (7.2,-1.7) {Paused};
\draw[inna,->] (l) -- node[al]{start} (c);
\draw[inna,->] (c) -- node[al]{3 s} (r);
\draw[inna,->] (r) -- node[al]{time over} (f);
\draw[inna,->] ([xshift=-5mm]r.south) -- node[al]{pause} ([xshift=-5mm]p.north);
\draw[inna,->] ([xshift=5mm]p.north) -- node[al]{resume} ([xshift=5mm]r.south);
\draw[inna,->] (f.north) -- ++(0,0.6) -| node[al, pos=0.25]{new game} (l.north);
\end{tikzpicture}
\caption{Phases of one game.}
\end{figure}

\textbf{Deliveries and points.} A player gets a package automatically when the motorcycle is inside
the area of the own team's store. M1 then picks a delivery point for this player. M4 gives a table
with the travel time from the store to every possible point. The times also count the obstacles that
are on the roads at that moment. M1 sorts this list and cuts it in the middle, so it gets a near
half and a far half. The admin sets a ``near chance'' for each player, and M1 uses it to choose the
half. A far point gives more points: points are a base value plus a bonus for each second of the
expected travel time. Both numbers are admin parameters. To deliver, the player presses
``Deliver''. The button is always active. M1 accepts it only inside the delivery radius and below
the speed limit. If M1 says no, it sends the reason, for example ``too fast'' or ``too far''. A fall
does not lose the package; the player only waits for the penalty time.

\begin{figure}[htbp]\centering
\begin{tikzpicture}[x=0.85cm, y=0.85cm]
\boardbase
\draw[tA, line width=2.2pt, rounded corners=3pt, opacity=0.9] (1,4) -- (1,3) -- (2,3) -- (2,2) -- (2.5,2);
\dpoint{2.5}{2}{tA}
\draw[tA!45, line width=1pt, dashed] (4,1.5) circle (0.18);
\bike{1}{3.5}{tA}{90}
\bike{3.5}{4}{tB}{0}
\bike{4}{2.5}{tD}{-90}
\bike{1.5}{1}{tC}{180}
\roadwork{3,2}{3,3}
\nmark{0.5}{5.2}{1}\nmark{0.55}{3.5}{2}\nmark{1.5}{2.7}{3}\nmark{2.5}{2.5}{4}\nmark{4.45}{1.5}{5}\nmark{3.35}{2.5}{6}
\node[anchor=west, font=\scriptsize, align=left] at (6.0,2.5) {\begin{tabular}{@{}c@{\ \ }l@{}}
\textbf{1} & Team store, two LEDs in team colour\\[2pt]
\textbf{2} & Player with a package, dot with a tail\\[2pt]
\textbf{3} & Fastest route found by M4\\[2pt]
\textbf{4} & Near delivery point, blinks in team colour\\[2pt]
\textbf{5} & A far point, worth more points\\[2pt]
\textbf{6} & Road work, so the route goes around\end{tabular}};
\end{tikzpicture}
\caption{One delivery as it will look on the LED board. The phone shows only a compass arrow.}
\end{figure}

\textbf{Obstacles and earthquakes.} M1 keeps one list of obstacles. Each obstacle has a type, a road
segment, a lane and a position. The list has a version number, and the number goes up on every
change. So M4 can see when it must compute routes again, and M5 and M8 can see when they must draw
again. All earthquakes take one path in M1, so a hand hit on the board, the motors and the admin
button act in the same way. The game strength is the measured strength times an admin factor. Roads
inside the earthquake radius get road work, and motorcycles inside the radius may fall.

\textbf{Admin parameters.} Every parameter has a name, a group, a minimum, a maximum and an owner
module. When the admin changes a value, M1 checks the limits first. Then it saves the value,
sends it to the owner (for example grip to M3 or traffic density to M4) and sends ``ok'' back to the
panel. A wrong value gets an error and the old value stays. Twice per second M1 also puts all graph
values (scores, delivery times, speeds, falls, traffic) into one metrics message.

\msection{Testing and demonstration}

M1 has no network code. So we can test every rule with fake M2, M3 and M4 parts that run inside the
test. The tests check the game flow, the lobby and the bots, the delivery cycle, the fair choice of
points, every admin command and the earthquake. One test plays a recorded game again with the same
seed and checks that the result is the same. In the demonstration, a full game with bots runs on the
main computer. The admin changes parameters live, starts an earthquake, and then plays the game log
again.

\end{document}
